\section{Protected cap constraints beyond added matchings}
\label{sec:path-six-protected-caps}

Matching augmentation is one special case of a fixed cap construction.
Branching tree caps give additional constraints on six-pole forest states.
These constraints exclude the earlier abstract half-integral obstruction.
They do not yet prove a bound for every bounded six-cut shore.

\begin{theorem}[A fixed cap cannot produce high terminal-path density]
\label{thm:path-six-fixed-cap-density}
Fix an integer $k\ge2$ and a finite loopless cap graph $F$.
Its vertices have the following types:
\begin{enumerate}
\item $q$ distinct attachment leaves of degree one.
\item $p$ designated terminal vertices of degree two.
\item Vertices of degree three.
\end{enumerate}
Every component of $F$ contains an attachment leaf.
Write $g=\beta(F)$ and $c=\operatorname{comp}(F)$.
Assume $q\le k$ and $k-q+p\ge2$.

Let $Y$ be a normalized $k$-pole of rank $h$ inside a cubic graph with no cyclic cut of size at most four.
Identify the attachment leaves of $F$ with any $q$ distinct ports of $Y$.
Call the resulting graph $Y_F$.
It is a normalized $(k-q+p)$-pole of rank
\[
 \beta(Y_F)=h+d_F,\qquad d_F=q+g-c.
\]
For every pair $a,b$ of its ports,
\[
 2^{-h}P_{ab}(Y_F)\le2^{d_F-1}+o(1)
 \qquad(h\to\infty).
\]
For a fixed finite collection of caps, the error is uniform over all caps and port choices.
\end{theorem}

\begin{proof}
The identifications make every filled port cubic.
All other degrees have the required values, and each cap component attaches to the connected graph $Y$.
Hence $Y_F$ is connected and normalized.
The rank increment is
\[
 |E(F)|-(|V(F)|-q)=g-c+q=d_F.
\]

Protect all new cap vertices and all filled ports of $Y$.
Their total number is fixed, independently of $h$.
If the normalized terminal-path density of $Y_F$ exceeded one half by a fixed amount along unbounded ranks, Theorem~\ref{thm:high-density-protected-cyclic-sets} would apply.
It would give two disjoint cyclic vertex sets avoiding every protected vertex and every remaining port.
Both sets lie in $Y$ and avoid all of its original ports.
Their induced graphs and boundaries are unchanged when the cap is removed and the original completion is restored.
Either set would therefore give a cyclic cut of size at most four in the original graph.
This contradiction proves the density bound.
Finiteness gives the stated uniformity.
\end{proof}

The path count in this theorem is a finite linear combination of the terminal forest states of $Y$.
Here is a precise coefficient description, including caps whose new vertices are the chosen terminals.
Let $A$ be the filled ports, and split the chosen pair into $T_Y\subseteq V(Y)$ and $T_F\subseteq V(F)\setminus A$.
For $S\subseteq A$, a path restricts to forests in $Y,F$ with odd sets $T_Y\cup S$ and $T_F\cup S$, respectively.
Their path pairings must join into one path between the chosen terminals when the vertices of $S$ are identified.
For each forest pairing in $Y$, its coefficient is the number of cap forests with this compatibility property.
The empty restriction is included when its odd set is empty.
These coefficients depend only on the fixed cap and terminal choices.

The description is exact.
A simple path restricts to two forests, since neither restriction can contain a cycle.
Conversely, compatible forests have maximum degree two after gluing, exactly two odd vertices, and one nonempty component.
Their union is therefore one simple terminal path.

\begin{corollary}[A three-port star cap]
\label{cor:path-six-star-cap}
Let $Y$ be a normalized six-pole of rank $h$ inside a cubic graph with no cyclic cut of size at most four.
Choose three ports $Q$ and a pair $e=ab$ among the other three.
For disjoint pairs $e,f$, let $B_{e|f}$ be the normalized sum of the two forest pairings other than $e|f$.
Then
\[
 a_e+\sum_{f\in\binom Q2}B_{e|f}\le2+o(1).
\]
There are sixty such inequalities.
\end{corollary}

\begin{proof}
Attach one new cubic vertex to the three ports of $Q$.
The resulting three-pole has rank $h+2$.
A terminal path either avoids the new vertex or uses exactly two of its incident edges.
Deleting those edges and the new vertex gives one of the two noncap forest pairings on $e\cup f$.
Every such forest conversely gives a terminal path.
Thus the displayed left side is exactly $2^{-h}$ times the augmented path count.
Apply Theorem~\ref{thm:path-six-fixed-cap-density}.
\end{proof}

\begin{corollary}[A four-leaf binary tree cap]
\label{cor:path-six-binary-tree-cap}
In the same setting, retain terminals $e=ab$ and partition the other four ports as $f\mid g$.
Let $R_{e|f|g}$ be the normalized six-terminal forest mass whose pairing has no edge in common with $\{e,f,g\}$.
Then
\[
 a_e+\sum_{f'\in\binom{[6]\setminus e}{2}}B_{e|f'}
     +R_{e|f|g}\le4+o(1).
\]
There are forty-five such inequalities.
\end{corollary}

\begin{proof}
Add two adjacent cubic vertices, attaching one to the two ports of $f$ and the other to the two ports of $g$.
The resulting two-pole has rank $h+3$.
For every two-port boundary pattern in the cap, its unique joining path gives one of the six terms $B_{e|f'}$.
For the four-port pattern, the cap's central edge is absent by parity.
Its two disjoint paths have pairing $f|g$.
They form one terminal path with a forest in $Y$ exactly for the six-terminal pairings counted by $R_{e|f|g}$.
The empty cap contributes $a_e$.
Theorem~\ref{thm:path-six-fixed-cap-density} now proves the inequality.
\end{proof}

These inequalities exclude the saved half-integral vector in Section~\ref{sec:six-cut-partial} as a limit of large poles in this setting.
For example, choose terminal pair $14$ and star attachment set $\{0,2,3\}$.
That vector gives
\[
 a_{14}+B_{14|02}+B_{14|03}+B_{14|23}=3,
\]
whereas Corollary~\ref{cor:path-six-star-cap} gives an upper limit of two.
The exact checker finds seven violated star inequalities and seven violated binary-tree inequalities for that vector.
This conclusion concerns the saved vector only.
It does not show that the enlarged relaxation has a strict universal deficit.
